\documentclass{beamer}

\usetheme{Berlin}
\usecolortheme{Orchid}
\useoutertheme{miniframes}
\setbeamertemplate{navigation symbols}{}

\usepackage[utf8]{inputenc}
\usepackage[T1]{fontenc}
\usepackage{amsmath}
\usepackage{amssymb}
\usepackage{booktabs}
\usepackage{graphicx}
\graphicspath{{../images/}}

\usepackage{xcolor}
\usepackage{listings}
\usepackage{hyperref}
\usepackage{tikz}
\usetikzlibrary{arrows.meta,positioning,shapes.geometric,calc}

\lstset{
  language=Java,
  basicstyle=\ttfamily\small,
  keywordstyle=\color{blue},
  commentstyle=\color{gray},
  stringstyle=\color{teal},
  showstringspaces=false,
  columns=fullflexible,
  breaklines=true
}

% Per-slide source line (references shown on the slide itself).
\newcommand{\slidesources}[1]{%
  \vfill
  {\tiny\color{gray}\raggedright\sloppy\parbox{\linewidth}{Sources: #1}\par}%
}

% Unit-wise source strings for this subject (used by slides/notes).
% Path is relative to the lecture `latex/` folder (the compilation CWD).
% OOPS with Java (CSEG1044) — unit-wise sources
%
% These are short, student-facing reference strings intended to be used with:
%   - \slidesources{...}  (in beamer slides)
%   - \notesources{...}   (in lecture notes)
%
% Style inspiration: other subjects in My-Subjects/ use semicolon-separated sources
% and include an "accessed" date for web documentation.

\newcommand{\OOPSUnitISources}{Schildt, \textit{Java: A Beginner's Guide} (10e), McGraw-Hill (Java fundamentals + OOP basics).; Oracle Java Tutorials: Getting Started + Language Basics + Classes and Objects (accessed \today).; Java SE API docs (e.g., \texttt{String}, \texttt{Scanner}) (accessed \today).}

\newcommand{\OOPSUnitIISources}{Schildt, \textit{Java: The Complete Reference} (12e), McGraw-Hill (classes, inheritance, packages, interfaces).; Bloch, \textit{Effective Java} (3e), Addison-Wesley, 2018 (design choices; interfaces vs abstract classes).; Oracle Java Tutorials: Classes and Objects + Interfaces + Packages (accessed \today).; Java SE API docs (e.g., \texttt{Comparable}, \texttt{Runnable}) (accessed \today).}

\newcommand{\OOPSUnitIIISources}{Schildt, \textit{Java: The Complete Reference} (12e) (nested/inner classes, exceptions, threads, I/O).; Oracle Java Tutorials: Nested Classes + Exceptions + Concurrency + Basic I/O (accessed \today).; Java SE API docs (e.g., \texttt{Thread}, \texttt{AutoCloseable}) (accessed \today).}

\newcommand{\OOPSUnitIVSources}{Schildt, \textit{Java: The Complete Reference} (12e) (generics, lambdas, Swing, JDBC).; Bloch, \textit{Effective Java} (3e) (generics + lambdas).; Oracle Java Tutorials: Generics + Lambda Expressions + Swing + JDBC Basics (accessed \today).; Java SE API docs (e.g., \texttt{java.util.function}, \texttt{javax.swing}, \texttt{java.sql}) (accessed \today).}

\newcommand{\OOPSUnitVSources}{Schildt, \textit{Java: The Complete Reference} (12e) (Collections Framework + wrapper classes).; Oracle Java docs: Collections Framework overview (accessed \today).; Java SE API docs (\texttt{java.util} collections; wrapper classes in \texttt{java.lang}) (accessed \today).}

\newcommand{\OOPSUnitVISources}{Git SCM: \textit{Pro Git} (2e) (version control workflow), accessed \today.; GitHub Docs (repositories, issues, pull requests), accessed \today.; Oracle Java Tutorials: Swing + JDBC (accessed \today).; JUnit 5 User Guide (accessed \today).; Apache Maven Project: Getting Started (accessed \today).; Gradle User Manual: Getting Started (accessed \today).; UML class diagrams: UML 2.x overview/spec (accessed \today).}

\newcommand{\OOPSMiscWhyInterfacesSources}{Schildt, \textit{Java: The Complete Reference} (12e) (interfaces).; Bloch, \textit{Effective Java} (3e) (prefer interfaces where appropriate).; Oracle Java Tutorials: Interfaces (accessed \today).; Java SE API docs (e.g., \texttt{Comparable}, \texttt{Runnable}) (accessed \today).}



\title[OOPS with Java]{Object Oriented Programming (CSEG1044)}
\subtitle{Unit 05 -- Lecture 02: Set/SortedSet + Map/SortedMap}
\begin{document}

\begin{frame}[fragile]
  \titlepage
  \vspace{0.5em}
  \slidesources{\OOPSUnitVSources}
\end{frame}

\begin{frame}{Quick Links}
  \centering
  \hyperlink{sec:overview}{\beamerbutton{Overview}}\hspace{0.6em}
  \hyperlink{sec:concepts}{\beamerbutton{Key Topics}}\hspace{0.6em}
  \hyperlink{sec:demo}{\beamerbutton{Demo}}\hspace{0.6em}
  \hyperlink{sec:summary}{\beamerbutton{Summary}}
  \slidesources{\OOPSUnitVSources}
\end{frame}

\begin{frame}{Agenda}
  \tableofcontents
  \slidesources{\OOPSUnitVSources}
\end{frame}

\section{Overview}
\label{sec:overview}

\begin{frame}{Lecture Snapshot}
\begin{itemize}[<+->]
  \item \textbf{Unit focus:} Collections Framework and Wrapper Classes
  \item \textbf{Course thread:} Use Set/SortedSet + Map/SortedMap to manage in-memory project data more cleanly and predictably inside Campus Resource Manager.
\end{itemize}
  \slidesources{\OOPSUnitVSources}
\end{frame}

\begin{frame}{Recall Before You Start}
\begin{itemize}[<+->]
  \item \textbf{List basics}: Sets and maps extend the collection story by optimizing for uniqueness and key-based lookup. Main reference: \texttt{Unit 05 / Lecture 01 slides/notes}.
  \item \textbf{Class design}: Set and map correctness depends on stable equality logic for stored objects. Main reference: \texttt{Unit 02 / Lecture 01 slides/notes}.
\end{itemize}
  \slidesources{\OOPSUnitVSources}
\end{frame}


\begin{frame}{Learning Outcomes}
  \begin{itemize}[<+->]
    \item Choose Set vs List and \texttt{HashSet} vs \texttt{TreeSet}.
    \item Explain \texttt{equals}/\texttt{hashCode} and common mistakes.
    \item Use \texttt{HashMap}/\texttt{TreeMap} for counting and sorted keys.
  \end{itemize}
  \slidesources{\OOPSUnitVSources}
\end{frame}

\begin{frame}{Key Terms to Know}
\small
\begin{columns}[t]
\column{0.48\textwidth}
\begin{itemize}
  \item \textbf{Set}: collection with no duplicates.
  \item \textbf{SortedSet}: set that maintains sorted order (e.g., \texttt{TreeSet}).
  \item \textbf{Hashing}: converting an object to an int hash value (\texttt{hashCode}).
\end{itemize}
\column{0.48\textwidth}
\begin{itemize}
  \item \textbf{Map}: key-value association.
  \item \textbf{SortedMap}: map with sorted keys (e.g., \texttt{TreeMap}).
  \item \textbf{Contract}: if \texttt{a.equals(b)} then \texttt{a.hashCode() == b.hashCode()} must hold.
\end{itemize}
\end{columns}
  \slidesources{\OOPSUnitVSources}
\end{frame}

\begin{frame}{Study Flow for This Lecture}
\begin{enumerate}[<+->]
  \item Refresh the prerequisite bridge before reading new ideas in isolation.
  \item Study the main build in this order: \textit{Set vs List; HashSet vs TreeSet} $\rightarrow$ \textit{Equality and hashing basics (must know)} $\rightarrow$ \textit{Map basics; HashMap vs TreeMap}.
  \item Trace the worked example and connect each line back to one concept frame.
  \item Attempt the mini exercises before moving to the recap and exit check.
\end{enumerate}
  \slidesources{\OOPSUnitVSources}
\end{frame}


\section{Key Topics}
\label{sec:concepts}

\begin{frame}{Unit 05: Collections and Wrapper Class}
  \begin{itemize}[<+->]
    \item Lecture focus: Set/SortedSet + Map/SortedMap
  \end{itemize}
  \slidesources{\OOPSUnitVSources}
\end{frame}

\begin{frame}{Today: Roadmap}
  \begin{itemize}[<+->]
    \item Set vs list; `HashSet` vs `TreeSet`
    \item Equality and hashing basics (`equals`/`hashCode`) and common mistakes
    \item Map basics; `HashMap` vs `TreeMap`; common use-cases (counting/grouping)
  \end{itemize}
  \slidesources{\OOPSUnitVSources}
\end{frame}

\begin{frame}{Set vs List}
  \begin{itemize}[<+->]
    \item List: ordered, duplicates allowed.
    \item Set: no duplicates; great for membership and deduplication.
  \end{itemize}
  \slidesources{\OOPSUnitVSources}
\end{frame}

\begin{frame}{HashSet vs TreeSet}
  \begin{itemize}[<+->]
    \item \texttt{HashSet}: no guaranteed order; fast average operations.
    \item \texttt{TreeSet}: sorted order; needs \texttt{Comparable} or \texttt{Comparator}.
  \end{itemize}
  \slidesources{\OOPSUnitVSources}
\end{frame}

\begin{frame}{Why \texttt{equals} and \texttt{hashCode} matter}
  \begin{itemize}[<+->]
    \item Hash-based collections use \texttt{hashCode} to find a bucket.
    \item \texttt{equals} confirms the match inside the bucket.
    \item Contract: if \texttt{a.equals(b)} then \texttt{a.hashCode()==b.hashCode()}.
  \end{itemize}
  \slidesources{\OOPSUnitVSources}
\end{frame}

\begin{frame}{Map basics}
  \begin{itemize}[<+->]
    \item Map = key $\rightarrow$ value.
    \item Use cases: counting, grouping, lookups by id.
    \item \texttt{HashMap}: fast average; \texttt{TreeMap}: keys sorted.
  \end{itemize}
  \slidesources{\OOPSUnitVSources}
\end{frame}

\begin{frame}{Checkpoint and Reflection}
\begin{block}{Reflection prompt}
Where would \textit{Set/SortedSet + Map/SortedMap} matter most in Campus Resource Manager, and which design choice would you justify using this lecture's ideas?
\end{block}
\begin{block}{Quick self-check}
Which part needs one more example before you move on: \textit{Set vs List; HashSet vs TreeSet}, \textit{Equality and hashing basics (must know)}, the worked example, or the recap?
\end{block}
  \slidesources{\OOPSUnitVSources}
\end{frame}

\section{Demo / Example}
\label{sec:demo}

\begin{frame}[fragile,shrink=8]{Example: Set + Map counting}
\begin{lstlisting}
import java.util.*;
public class Demo {
    public static void main(String[] args) {
        List<String> names = Arrays.asList("Asha", "Ravi", "Asha");
        Set<String> unique = new HashSet<>(names);
        System.out.println(unique);

        String[] words = {"a","b","a","c","b","a"};
        Map<String,Integer> freq = new HashMap<>();
        for (String w : words) freq.put(w, freq.getOrDefault(w,0) + 1);
        System.out.println(freq);

        Map<String,Integer> sorted = new TreeMap<>(freq);
        System.out.println(sorted);
    }
}
\end{lstlisting}
  \slidesources{\OOPSUnitVSources}
\end{frame}

\begin{frame}{Mini Exercises (2--3 mins)}
  \begin{enumerate}[<+->]
    \item When do you prefer Set over List?
    \item Which keeps sorted order: \texttt{HashSet} or \texttt{TreeSet}?
    \item If you override \texttt{equals}, what else must you override?
  \end{enumerate}
  \slidesources{\OOPSUnitVSources}
\end{frame}

\begin{frame}{Watch-outs}
\begin{itemize}[<+->]
  \item Forgetting the equals/hashCode contract and blaming the collection.
  \item Expecting HashSet or HashMap to preserve sorted order.
  \item Using a map or set before confirming the problem truly needs keys or uniqueness.
\end{itemize}
  \slidesources{\OOPSUnitVSources}
\end{frame}

\section{Summary}
\label{sec:summary}

\begin{frame}{Key Takeaways}
  \begin{itemize}[<+->]
    \item Set enforces uniqueness; TreeSet provides sorted order.
    \item Hashing relies on correct \texttt{equals}/\texttt{hashCode}.
    \item Map enables fast key-based lookup and counting/grouping.
  \end{itemize}
  \slidesources{\OOPSUnitVSources}
\end{frame}

\begin{frame}{Checkpoint Questions}
  \begin{enumerate}[<+->]
    \item Why do Hash-based collections need \texttt{hashCode} and \texttt{equals}?
    \item When do you choose \texttt{TreeMap} over \texttt{HashMap}?
    \item Give one example of counting/grouping using a map.
  \end{enumerate}
  \slidesources{\OOPSUnitVSources}
\end{frame}

\begin{frame}{Next Study Steps}
\begin{itemize}[<+->]
  \item Revisit the recall references if any older idea still feels weak.
  \item Rework the lecture example without looking at the notes first.
  \item Attempt the self-study practice and then answer the exit questions in full sentences.
  \item Apply the idea once inside Campus Resource Manager to make the concept stick.
\end{itemize}
  \slidesources{\OOPSUnitVSources}
\end{frame}

\begin{frame}{Next Lecture Preview}
  \textbf{Next:} Wrapped Collections + \texttt{Collections} utilities + Wrapper classes
  \vspace{0.6em}
  \begin{itemize}[<+->]
    \item Utility methods and common pitfalls with autoboxing.
  \end{itemize}
  \slidesources{\OOPSUnitVSources}
\end{frame}

\end{document}
